% uid: ""
% school: "lycee"
% level: "1ere"
% chapter: "Activités numériques"
% title: "Chiffres cachés"
% difficulty: "Moyen"
% tags: ["divisibilite", "criteres-de-divisibilite", "chiffres-caches", "multiples-de-75"]
% provenance: "own"
% batch: "2026-11-22-revision-1ere-A"

% ---STATEMENT---
Dans chaque cas, déterminer tous les chiffres $a$ et $b$ tels que :
\begin{enumerate}
  \item le nombre $\overline{2a56}$ soit divisible par $8$ ;
  \item le nombre $\overline{7a4b}$ soit divisible par $8$ et par $9$ ;
  \item le nombre $\overline{a3b5}$ soit un multiple de $75$ (avec $a\neq0$).
\end{enumerate}

% ---SOLUTION---
\begin{prop}[Critères de divisibilité]
Un entier est divisible par \textbf{8} si le nombre formé par ses \ib{trois derniers chiffres} l'est, par \textbf{25} si ses \ib{deux derniers chiffres} forment $00$, $25$, $50$ ou $75$, par \textbf{3} (ou \textbf{9}) si la \ib{somme de ses chiffres} l'est.

Si deux entiers sont \textbf{premiers entre eux} et divisent $n$, alors leur produit divise $n$.
\end{prop}
\begin{enumerate}
  \item Le critère de $8$ porte sur $\overline{a56}=100a+56$.

  On a $56=7\times8$ : il reste à avoir $8\mid100a$.

  On écrit $100a=96a+4a$, avec $96a=8\times12a$ :
  \[8\mid100a\iff8\mid4a\iff a\ \text{est pair}\]
  Les chiffres possibles sont donc \ib{$a\in\{0,2,4,6,8\}$}.

  Contrôle : $056=7\times8$, $256=32\times8$, $456=57\times8$, $656=82\times8$, $856=107\times8$.

  \item \textbf{Condition pour $8$.} On regarde $\overline{a4b}=100a+40+b$.
  \[100a+40+b=8\,(12a+5)+4a+b\]
  Donc $8\mid\overline{7a4b}\iff8\mid4a+b$.

  \textbf{Condition pour $9$.} La somme des chiffres est $7+a+4+b=11+a+b$, avec $11\leq11+a+b\leq29$.

  Elle est multiple de $9$ si elle vaut $18$ ou $27$ :
  \[a+b=7\qquad\text{ou}\qquad a+b=16\]
  On écrit $4a+b=3a+(a+b)$ pour chaque cas.

  \textbf{Cas $a+b=7$.} On veut $8\mid3a+7$ avec $0\leq a\leq7$ :
  \[3a+7\in\{7,10,13,16,19,22,25,28\}\]
  Seul $16$ est un multiple de $8$ : $a=3$, donc $b=4$.

  \textbf{Cas $a+b=16$.} On veut $8\mid3a+16$, c'est-à-dire $8\mid3a$, donc $8\mid a$ (car $3$ et $8$ sont premiers entre eux).

  Avec $7\leq a\leq9$ (pour que $b=16-a\leq9$), on obtient $a=8$ et $b=8$.

  Les nombres sont \ib{$7344$} et \ib{$7848$}.

  Contrôle : $7344=8\times918=9\times816$ et $7848=8\times981=9\times872$.

  \item On a $75=3\times25$, avec $3$ et $25$ premiers entre eux : $\overline{a3b5}$ est multiple de $75$ si et seulement s'il est divisible par $25$ \textbf{et} par $3$.

  \textbf{Condition pour $25$.} Les deux derniers chiffres $\overline{b5}$ doivent valoir $25$ ou $75$ :
  \[b=2\qquad\text{ou}\qquad b=7\]
  \textbf{Condition pour $3$.} La somme des chiffres est $a+3+b+5=a+b+8$.

  \textbf{Si $b=2$ :} $a+10$ est multiple de $3$, avec $1\leq a\leq9$ : $a\in\{2,5,8\}$.

  \textbf{Si $b=7$ :} $a+15$ est multiple de $3$, donc $a$ est multiple de $3$ : $a\in\{3,6,9\}$.

  On obtient six nombres :
  \[\ib{$2325,\ 5325,\ 8325,\ 3375,\ 6375,\ 9375$}\]
  Contrôle : $2325=31\times75$, $3375=45\times75$, $9375=125\times75$.
\end{enumerate}
\textbf{Idée :} ramener chaque critère à une condition sur les chiffres inconnus, puis croiser les conditions.

\textbf{Piège :} oublier le deuxième multiple de $25$ ($25$ et $75$) ou accepter $a=0$ pour un chiffre de tête.
